% ============================================================
% CHAPTER 2. ANALYSIS OF BEST PRACTICES AND ACADEMIC LITERATURE
% ============================================================

\section{Analysis of Best Practices and Academic Literature}
\label{sec:literature}

This chapter positions the project within two connected bodies of knowledge. The first is academic research on financial machine learning, chart-image forecasting, intraday market prediction, time-series representation and explainable artificial intelligence. The second is the practical decision-support tradition in which analytical systems are designed not only to produce predictions, but also to organize evidence, support user judgment, and control the risk of over-reliance on automated outputs.

The review is structured as a taxonomy rather than as a chronological list of papers, following the logic of systematic review guidelines while remaining focused on the project problem \cite{kitchenham_slr,prisma2020}. This is important for the project-based nature of the thesis. The practical problem is not simply to train a CNN. The practical problem is to design a selective, auditable decision-support artifact for intraday trading analysis. Therefore, the literature is reviewed through six themes: financial machine learning and chart-image prediction; intraday and high-frequency financial deep learning; time-series-to-image learning and visual feature baselines; explainable AI and saliency-map limitations; decision-support systems and dashboard artifacts; and the final research gap.

\subsection{Financial Machine Learning and Chart-Image Prediction}
\label{subsec:financial_ml_chart_images}

\begin{figure*}[!t]
\centering
\includegraphics[width=\textwidth]{Figure_2_reimag_alignment.pdf}
\caption{Methodological alignment between Jiang, Kelly and Xiu's Re-Imag chart-image design and the present intraday MOEX Si futures decision-support study. Re-Imag uses daily equity chart images and portfolio-style evaluation, while this study uses intraday futures images and selective dashboard-based decision support.}
\label{fig:reimag-alignment}
\end{figure*}



Financial forecasting is a difficult domain for machine learning because the signal-to-noise ratio is low, labels are unstable, and evaluation can be distorted by leakage, non-stationarity and selection bias. Modern financial machine learning therefore emphasizes careful data handling, chronological validation and skepticism toward purely in-sample performance \cite{lopez_de_prado_afml,bailey_deflated_sharpe,bailey_probabilistic_sharpe}. Deep learning has expanded the model families available for financial forecasting, including supervised forecasting and reinforcement-learning approaches to trading and portfolio management, but the literature continues to stress that predictive performance should be interpreted cautiously, especially when the task is short-horizon return direction \cite{heaton_deep_learning_finance,sezer_financial_forecasting_survey,lim_zohren_time_series_forecasting,jiang_portfolio_drl,moody_direct_reinforcement}.

A central strand of the literature studies whether visual representations of prices contain predictive information. Traditional technical analysis assumes that historical price and volume patterns may contain recurring structures. Earlier empirical studies of technical trading rules and chart patterns provide a foundation for treating price shapes as objects of analysis \cite{brock_technical_trading,lo_foundations_technical_analysis}. However, these studies typically rely on predefined rules or handcrafted pattern definitions. Machine learning changes this setting by allowing representations to be learned from data rather than specified entirely by the researcher.

The primary methodological benchmark for this thesis is Jiang, Kelly and Xiu's ``(Re-)Imag(in)ing Price Trends'' \cite{jiang_reimagining_price_trends}. Their study converts stock price histories into chart images and trains CNNs to forecast return direction. The paper is central to the present thesis for four reasons. First, it provides a rigorous justification for treating financial price histories as images. Second, it demonstrates that CNNs can learn visual price-pattern information from OHLC and volume-style chart representations. Third, it evaluates multiple visual horizons, including 5-, 20- and 60-day image windows, which motivates the multi-scale extension in the present project. Fourth, it provides a benchmark for comparing CNN image learning with more traditional financial predictors and visual feature alternatives.

This thesis follows the core idea of Re-Imag but adapts it to a different applied problem. The benchmark study operates in a daily equity setting, while the present study works with fifteen-minute MOEX Si futures data. The benchmark study is oriented toward return prediction and portfolio evaluation, while this thesis is oriented toward selective decision support and analyst inspection. The difference is not incidental: it defines the contribution. The present project asks whether chart-image CNN logic can be transformed into an intraday futures decision-support artifact that abstains from ambiguous cases and links selected cases to raw chart and XAI evidence.

Financial time-series forecasting surveys provide further context. Sezer, Gudelek and Ozbayoglu review deep learning approaches in financial forecasting and show the importance of CNNs, recurrent models and hybrid architectures in price-prediction applications \cite{sezer_financial_forecasting_survey}. Lim and Zohren review deep learning for time-series forecasting more broadly, emphasizing representation learning, sequence modeling, interpretability and uncertainty \cite{lim_zohren_time_series_forecasting}. These surveys support the methodological choice to use deep visual models, but they also reinforce that forecasting systems should be evaluated in relation to horizon, data regime and decision use case rather than by one accuracy score alone.

The literature therefore suggests a cautious principle: chart-image CNNs are a legitimate research direction, but their practical value depends on the decision context, validation design and interpretation layer. This principle directly motivates the selective and explainable design of the present system.

\subsection{Intraday and High-Frequency Financial Deep Learning}
\label{subsec:intraday_hft}

Intraday financial prediction differs from daily prediction because short-horizon returns are affected by microstructure effects, liquidity, order flow, intraday seasonality and market-impact considerations. Market microstructure theory emphasizes that observed prices reflect the interaction of information, liquidity provision and trading mechanisms \cite{ohara_microstructure}. For this reason, intraday forecasting models often rely on more detailed representations of market dynamics than daily models.

Deep learning has been widely applied to high-frequency and limit-order-book data. DeepLOB is a prominent example, using deep neural networks to forecast price movement from limit-order-book states \cite{deeplob}. Other studies use CNNs and recurrent models to learn structure from order-book depth, order-flow dynamics and high-frequency price formation \cite{zhang_lob_cnn,tsantekidis_lob,sirignano_cont_universal}. Support-vector and other machine-learning approaches have also been used to model limit-order-book dynamics \cite{kercheval_lob_svm}. These studies show that deep models can extract structure from high-dimensional intraday market data, but they also highlight the complexity of evaluation under market microstructure noise.

The present thesis differs from limit-order-book prediction because it does not use order-book depth or tick-level order flow. Instead, it uses fifteen-minute OHLC chart images constructed from Si futures candles. This is a coarser representation, but it is more directly connected to trader-facing chart interpretation and dashboard-based visual audit. The model therefore sits between daily chart-image learning and high-frequency microstructure modeling. It is intraday, but not order-book based.

This intermediate position has practical advantages and limitations. The advantage is interpretability and operational simplicity: chart images can be inspected by an analyst and linked to familiar candlestick or price-path structures. The limitation is that some microstructure information is lost. The value/turnover field is unavailable in the supplied data because it remains all zero, so the final system cannot claim to exploit valid turnover or transaction-value information. This limitation strengthens the case for caution and for the use of abstention rather than universal prediction.

The intraday literature also supports the need for chronological validation. Random train/test splitting would be inappropriate because adjacent windows overlap and future information could leak into training. The present project therefore uses chronological data organization and treats overlapping images as a source of dependence that must be discussed, not ignored.

\subsection{Time-Series-to-Image Learning and Visual Feature Baselines}
\label{subsec:tsi_visual_baselines}

A broader time-series representation literature supports the conversion of temporal data into images. Wang and Oates proposed encoding time series as images, showing that visual representations can enable CNN-based classification of temporal patterns \cite{wang_oates_imaging_time_series}. Related representations include Gramian Angular Fields, Markov Transition Fields and recurrence plots, all of which transform temporal relationships into two-dimensional structures that can be processed by computer-vision models \cite{recurrence_plots,gramian_fields}. In financial applications, Sezer and Ozbayoglu show how time-series-to-image approaches can be used for trading-related prediction tasks \cite{sezer_ozbayoglu_tsi}.

Deep learning for time-series classification is also relevant. Fawaz et al. review deep neural networks for time-series classification and show that CNN-style architectures are a major family of models for learning temporal patterns \cite{fawaz_tsc_review}. InceptionTime and ROCKET represent two important directions in time-series classification: deep convolutional ensembles and large randomized convolutional feature transformations \cite{inceptiontime,rocket}. The hctsa framework and shapelet methods represent another family of approaches based on large-scale feature extraction or discriminative subsequences \cite{hctsa,shapelets,learning_shapelets,random_dilated_shapelets}. Distance-based methods such as dynamic time warping and DBA provide additional baselines for comparing temporal shapes \cite{dtw,dba}.

The present thesis does not attempt to benchmark every time-series classification method, nor does it use conditional visual-reasoning layers such as FiLM; those architectures remain adjacent background for future model extensions \cite{perez_film}. Instead, it uses the literature to justify the representation choice and to interpret visual baselines. The chart image is a direct visual representation, not a transformed recurrence or Gramian field. It is selected because the project is explicitly trader-facing: analysts can inspect candlestick-style chart images more naturally than abstract time-series images.

Classical computer-vision features are used as important comparators. HOG captures distributions of local gradient orientations and is well suited for edge and shape information \cite{dalal_hog}. HAAR-like features summarize contrast patterns over rectangular regions and have a long history in computer vision \cite{viola_haar}. In the Re-Imag context, such classical visual features are relevant because they test whether predictive information can be extracted from chart geometry without a deep CNN. In this thesis, the HOG tail baseline performs meaningfully, reaching 0.704 accuracy, 0.754 balanced accuracy and 0.455 MCC at 2.0\% coverage. HAAR-like features are weaker, with the best reported row reaching 0.574 accuracy, 0.582 balanced accuracy and 0.155 MCC at 5.1\% coverage. This comparison supports a balanced conclusion: visual structure matters, but the strongest selective CNN result remains superior in the supplied experiment package.

\subsection{Explainable AI and Saliency-Map Limitations}
\label{subsec:xai_literature}

Explainability is essential in this thesis because the proposed artifact is intended for human-in-the-loop decision support. A black-box score alone is insufficient for analyst trust, especially in a financial domain where model errors can have operational consequences. General XAI surveys define the field and distinguish transparent models, post-hoc explanations, local explanations and global explanations \cite{guidotti_survey,adadi_berrada_xai,arrieta_xai}. XAI in finance is particularly relevant because financial applications require accountability, risk awareness and practical interpretability \cite{weber_xai_finance}.

For deep neural networks, visual explanation methods are a major class of post-hoc interpretation. Class Activation Mapping (CAM) and Grad-CAM highlight image regions associated with class-discriminative decisions \cite{zhou_cam,selvaraju_gradcam}. Grad-CAM++ refines Grad-CAM for improved localization in some settings \cite{chattopadhay_gradcampp}. Score-CAM, SmoothGrad, Layer-wise Relevance Propagation, Integrated Gradients, DeepLIFT, LIME, SHAP, counterfactual explanations, influence functions and TCAV represent related or alternative approaches for interpreting black-box models \cite{wang_scorecam,smilkov_smoothgrad,montavon_lrp,sundararajan_integrated_gradients,shrikumar_deeplift,ribeiro_lime,lundberg_shap,wachter_counterfactual,koh_influence,kim_tcav}. These methods vary in assumptions, computational cost and explanation target.

However, the XAI literature also warns against overclaiming saliency maps. Adebayo et al. show that some saliency methods can fail sanity checks, meaning that visually plausible explanations do not always reflect learned model parameters \cite{adebayo_sanity_checks}. Lipton cautions that interpretability is often used imprecisely, and Rudin argues that in high-stakes applications, interpretable models should be preferred when possible rather than relying uncritically on explanations of black boxes \cite{lipton_mythos,rudin_black_boxes}. Miller emphasizes that human explanations are selective, contrastive and social, not merely technical heatmaps \cite{miller_explanation}. These works are critical for this thesis because they justify a cautious interpretation of Grad-CAM.

The present project uses Grad-CAM as an audit layer, not as proof of causal reasoning. Initial analysis of heatmap-only images suggested broad or possibly background-heavy activation. Later masked-overlay, full-overlay and raw selected-case images refined the interpretation: in representative selected cases, activation is generally chart-anchored but spatially diffuse. This means the model is not simply focusing on empty background, but the highlighted regions extend beyond the precise candle path. Accordingly, Grad-CAM is interpreted as support for pattern-level inspection and diagnostic review, not as evidence of exact candle-level causal explanation.

A closely related applied precedent is the KOSPI 200 futures CNN and Grad-CAM study, which is useful because it connects futures prediction, CNNs, chart-style data and Grad-CAM-style visualization \cite{kim_kospi_gradcam}. The present thesis differs by emphasizing selective abstention, dashboard workflow, and MOEX Si futures decision support.

\subsection{Decision-Support Systems, Human-in-the-Loop Analytics and Dashboard Artifacts}
\label{subsec:decision_support}

The project-based contribution of this thesis is not only a model, but an artifact. Decision-support systems translate data and models into structured assistance for human decisions. In this thesis, the Decision Support Studio dashboard performs that translation by organizing selected signals, displaying model confidence, showing visual evidence, and making limitations visible.

Design science research provides an appropriate framework for this artifact orientation. Hevner et al. describe design science as the creation and evaluation of artifacts intended to solve relevant organizational problems \cite{hevner_dsr}. Peffers et al. provide a design science research methodology that includes problem identification, objectives, design and development, demonstration, evaluation and communication \cite{peffers_dsrm}. FEDS emphasizes that artifact evaluation should be aligned with the purpose and context of the designed system \cite{venable_feds}. These frameworks support the thesis's project-based structure: the dashboard is evaluated not as a commercial product but as a research artifact that operationalizes selective intraday decision support.

Human-in-the-loop analytics are particularly important when model predictions are uncertain and explanations are imperfect. A responsible dashboard should not hide uncertainty. It should show coverage, sample size, confidence, baseline comparison and explanation quality. The Decision Support Studio is designed around this principle. It includes KPI cards, a signal queue, case-level image views, benchmark ladder, selective coverage frontier, XAI diagnostics and governance notes. These components make the artifact practical and academically defensible.

This decision-support framing also protects the thesis from overclaiming. The dashboard does not present the CNN as a fully automated trading agent. Instead, it supports analyst triage. A surfaced case is an object for review, not a command to trade. This is consistent with the empirical results: the full-sample model edge is modest, while selected tail cases are stronger but limited in coverage.

\subsection{Research Gap and Positioning}
\label{subsec:research_gap}

The reviewed literature establishes five points. First, financial machine learning can extract patterns from market data, but evaluation must be cautious because of noise, leakage and overfitting risk. Second, chart-image CNNs are a strong and relevant methodological direction, especially after ``(Re-)Imag(in)ing Price Trends.'' Third, intraday forecasting is more difficult than daily chart prediction because of microstructure noise, overlapping windows and short-horizon instability. Fourth, XAI can support audit and communication but cannot prove causal reasoning by itself. Fifth, project-based decision-support artifacts require more than model accuracy: they require workflow, governance and human inspection.

The gap addressed by this thesis is therefore the following:

\begin{quote}
Prior work shows that chart-image CNNs can forecast returns in equity markets, but less is known about selective, explainable, human-in-the-loop intraday futures decision support using multi-scale chart images, Grad-CAM audit and dashboard-based analyst workflow.
\end{quote}

The present study addresses this gap by adapting the chart-image CNN logic of Jiang, Kelly and Xiu to MOEX Si futures, extending it with multi-scale intraday input windows, comparing it with numeric and visual baselines, using selective probability-tail rules to abstain from ambiguous windows, auditing selected cases with Grad-CAM, and implementing the Decision Support Studio dashboard as a practical decision-support artifact.

The resulting contribution is deliberately bounded. The thesis does not claim to prove market inefficiency, create an automated profitable trading system or establish causal candlestick reasoning. It claims to design and evaluate a selective, auditable and human-in-the-loop decision-support system for intraday futures analysis.
